\scanpage{9}
% Pre-ruled music paper.
\uline{Cor} 1 \deleted{\unsure{Il en sera}} \deleted{\unsure{Dans ces conditions}} On obtient encore des conditions équivalentes en exigeant \inserted{de plus} sous a) que les $L_n$, \deleted{$\pi_n(L_{\bullet})$} soient de type fini, en a') que les $L_n^!$, \deleted{$H_n(L_{\bullet}^!)$} soient de type fini, dans b) et b') que les $M_i$ et $N_j$ soient de type fini.

On \unsure{aura|a alors} dans les conditions de prop 1
\[
 L_n=\prod_{1\leq i\leq n+1}\Lambda^{i}\Phi_{n k}\otimes_k M_i
 \times\prod_{\unsure{0|1}\leq j\leq n\deleted{\unsure{+1}}}\Lambda^j\Psi_n\otimes N_j
\]
\deleted{d'où on trouve}

\uline{Cor} 2 Sous les conditions équivalentes \deleted{de} \deleted{la} \deleted{\unsure{la prop., pour que $L_n$ soit de type fini, pour}} \inserted{\unsure{précédentes}, pour que $L_{\bullet}^!$ soit} à degrés bornés, il faut\unsure{)}, qu'il n'y ait qu'un nombre \inserted{fini} d'indices $i$ et $j$ tels que $M_i\ne0$, $N_j\ne0$.

\deleted{Sous les conditions} On appellera \unsure{les|un} $L_{\bullet}$ \deleted{\unsure{strict}} \deleted{\unsure{parfait}} \deleted{\unsure{spécial}} \inserted{\uline{\unsure{spécial}}} si $L_{\bullet}^!$ est \deleted{str. parfait, i.e.} \inserted{à degrés bornés, à composantes} et à $H_n$ projectifs de t.f. --- i.e. si $L_{\bullet}$ est \inserted{$\simeq$} somme finie de modules de la forme \deleted{\unsure{$k$}} $\Lambda^{\unsure{j|i}}\Phi_{*k}\otimes_k M_i$ \inserted{$(i\geq1)$}, $\Lambda^j\Psi_*\otimes M_j$ $(j\geq0)$, avec les $M_i$ et $N_j$ proj. de t.f., \uline{strictement spéciaux} si de plus les $M_i,N_j$ sont \uline{libres}. Il revient au $\overline{\mathrm m}$ de
% archive sheet number: 8
